\documentclass[11pt]{article}
\usepackage[margin=1in]{geometry}
\usepackage{amsmath,amssymb,amsthm,mathtools}
\usepackage[hidelinks]{hyperref}
\usepackage{microtype}
\usepackage{enumitem}

\newtheorem{theorem}{Theorem}[section]
\newtheorem{proposition}[theorem]{Proposition}
\newtheorem{lemma}[theorem]{Lemma}
\newtheorem{corollary}[theorem]{Corollary}
\theoremstyle{definition}
\newtheorem{definition}[theorem]{Definition}
\theoremstyle{remark}
\newtheorem{remark}[theorem]{Remark}

\newcommand{\N}{\mathbb N}
\newcommand{\Col}{\operatorname{Col}}
\newcommand{\Reach}{\rightsquigarrow}
\newcommand{\card}[1]{\lvert #1\rvert}

\title{Quantitative Remainder Transport in Collatz Dynamics}
\author{Shreyas Samir Kolte}
\date{November 2026}

\begin{document}
\maketitle

\begin{abstract}
Let $\Col(n)=n/2$ for even $n$ and $\Col(n)=3n+1$ for odd $n$.
We combine an explicit positive-density theorem for Collatz convergence in
logarithmic time with effective residue transport in inverse Collatz trees.
For every modulus $N\ge 1$, every residue $r\pmod N$, and every sufficiently
large parameter $X$, we construct a finite set $W_{N,r}(X)$ of distinct
integers such that every $n\in W_{N,r}(X)$ satisfies
\[
 n\equiv r\pmod N,\qquad n\Reach 1,
\]
and
\[
 \card{W_{N,r}(X)}\ge
 \frac{\delta X}{2(\lceil (523/50)\log X\rceil+4)}.
\]
Here $\delta>0$ and the cutoff are the explicit constants in Mazur's
positive-density formalization.  The general construction gives
\[
 n\le (3X+1)2^{2^{48N^4}},
\]
and we prove substantially sharper bounds when $N$ is coprime to $6$, when
$N$ is a power of $3$, and when a factorization certificate for $N$ and the
target residue is supplied.  All principal statements are formalized in Lean
4.30.0-rc2 and their axiom reports contain only \texttt{propext},
\texttt{Classical.choice}, and \texttt{Quot.sound}.  The result does not prove
the Collatz conjecture or positive density inside each residue class; it
isolates orbit coalescence as the source of the logarithmic loss.
\end{abstract}

\section{Introduction}
The Collatz map is the function $\Col:\N_{>0}\to\N_{>0}$ defined by
\[
 \Col(n)=\begin{cases}n/2,&2\mid n,\\3n+1,&2\nmid n.\end{cases}
\]
The Collatz conjecture asserts that every positive integer eventually reaches
$1$.  Although the rule is elementary, its forward and inverse dynamics
combine deterministic arithmetic with behavior that repeatedly resembles a
mixing process.  This resemblance is useful but dangerous: qualitative
coverage of residue classes is much weaker than equidistribution, and neither
implies universal convergence.

This paper studies a precise interface between two rigorous phenomena.  The
first is abundance: a fixed explicit positive proportion of starting values
reaches $1$ within $(523/50)\log n$ ordinary Collatz steps.  The second is
transport: for a target not divisible by $3$, a predecessor can be found in
every prescribed residue class, with an explicit upper bound.  Applying
transport separately to many fast-convergent targets does not automatically
produce equally many distinct predecessors, because several targets may lie
on one forward orbit.  We prove that the multiplicity is at most logarithmic,
which yields an explicit $X/\log X$ family in every residue class.

The central conclusion is therefore stronger than the qualitative assertion
that every residue occurs infinitely often in an inverse tree.  It gives a
finite, explicit, growing population of integers in each residue class that
are certified to reach $1$.

\section{Prior work}
Terras proved fundamental density results for finite stopping time, with
related work by Everett.  Wirsching developed a systematic analysis of the
inverse Collatz graph and later formulated the positive-predecessor-density
problem.  Tao proved that almost all Collatz orbits attain any prescribed
diverging bound, in logarithmic density, by an approximate transport argument
for the Syracuse map.

Mazur subsequently gave an explicit positive-density theorem in logarithmic
ordinary time: there are explicit constants $\delta>0$ and $X_0$ such that,
for every $X\ge X_0$, at least $\delta X$ positive integers below $X$ reach
$1$ within $(523/50)\log n$ steps.  A separate predecessor-density result
shows positive lower density for the predecessors of each fixed target not
divisible by $3$.

Hallqvist's Every Remainder Everywhere theorem proves that the inverse tree of
every target not divisible by $3$ meets every residue class modulo every
modulus infinitely often.  Its effective version supplies the explicit
general witness bound used here.  These results address modular accessibility,
not the number of distinct transported witnesses obtained from a population
of targets.

Lau and AxiomProver recently made positive-density constants more explicit and
proved target-specific predecessor results.  Those results concern abundance
for a prescribed target.  The present theorem addresses a different axis:
simultaneous quantitative transport into each prescribed congruence class.
No claim of priority beyond the exact residue-wise statements formalized here
is made.

\section{Definitions}
Write $a\Reach b$ when an iterate of $\Col$ starting at $a$ equals $b$.
For $C\ge0$ and $X\in\N$, let
\[
 \mathcal F_C(X)=\{n<X:n\text{ is odd and reaches }1
       \text{ in at most }C\log n\text{ steps}\}.
\]
Define the live-target map
\[
 L(n)=\begin{cases}3n+1,&3\mid n,\\n,&3\nmid n.\end{cases}
\]
For odd $n$, either $L(n)=n$ or $L(n)=\Col(n)$.  Thus $L(n)$ is not divisible
by $3$, reaches $1$, and requires no more steps to reach $1$ than $n$ does.
Let
\[
 \mathcal T_C(X)=L(\mathcal F_C(X)).
\]
The map $L$ has fibres of size at most two on odd inputs, hence
\begin{equation}\label{eq:live-count}
 \card{\mathcal F_C(X)}\le2\card{\mathcal T_C(X)}.
\end{equation}

\section{The orbit-fibre bound}
Suppose a selector $s$ assigns to each $w\in\mathcal T_C(X)$ a predecessor
$s(w)$ with $s(w)\Reach w$.  If $s(u)=s(v)$, then both $u$ and $v$ occur on
the same forward orbit, namely the orbit of the common selected predecessor.
If every target reaches $1$ within $B$ steps, all such targets lie among a
prefix of length $B+1$ together with the terminal cycle $1,4,2$.  Consequently
each fibre of $s$ has cardinality at most $B+4$.

\begin{proposition}[Selected-witness counting]
If every $w$ in a finite target set $S$ reaches $1$ within $B$ steps and
$s(w)\Reach w$, then
\[
 \card S\le \card{s(S)}(B+4).
\]
\end{proposition}

For $C=523/50$, one may take
\[
 B=B(X)=\left\lceil\frac{523}{50}\log X\right\rceil.
\]
Combining the proposition with \eqref{eq:live-count} gives
\begin{equation}\label{eq:collision}
 \card{\mathcal F_{523/50}(X)}
 \le2\card{s(\mathcal T_{523/50}(X))}(B(X)+4).
\end{equation}

\section{Effective residue transport}
Effective Remainder Everywhere supplies, for every $N\ge1$, residue $r$, and
live target $w$, a selected predecessor $s_{N,r}(w)$ such that
\[
 s_{N,r}(w)\Reach w,\qquad s_{N,r}(w)\equiv r\pmod N,
\]
with
\begin{equation}\label{eq:generic-bound}
 s_{N,r}(w)\le w2^{2^{48N^4}}.
\end{equation}
Since every target in $\mathcal T_C(X)$ is below $3X+1$, the right-hand side
of \eqref{eq:generic-bound} is at most
$(3X+1)2^{2^{48N^4}}$.

\section{Main theorems}
Let $\delta>0$ and $X_0$ be Mazur's explicit density constant and cutoff.

\begin{theorem}[Quantitative direct-to-one remainder transport]
For every $N\ge1$, $r\in\N$, and $X\ge X_0$, there is a finite set
$W_{N,r}(X)$ satisfying
\[
 \card{W_{N,r}(X)}\ge
 \frac{\delta X}{2(B(X)+4)}.
\]
Every $m\in W_{N,r}(X)$ satisfies
\[
 m\equiv r\pmod N,\qquad m\Reach1,
 \qquad m\le(3X+1)2^{2^{48N^4}}.
\]
\end{theorem}

\begin{proof}
Mazur's theorem gives
$\delta X\le\card{\mathcal F_{523/50}(X)}$.  Apply the effective residue
selector to every live target and then use \eqref{eq:collision}.  The selected
image is a finite set of distinct integers in the prescribed residue.  Each
selected integer reaches its target and each target reaches $1$, so
transitivity gives convergence to $1$.  The numerical bound follows from
\eqref{eq:generic-bound} and $w<3X+1$.
\end{proof}

\begin{theorem}[Coprime-to-six refinement]
If $\gcd(N,6)=1$, the same cardinality conclusion holds with
\[
 m\le(3X+1)2^{2^{15(3N^2)}}
       =(3X+1)2^{2^{45N^2}}.
\]
\end{theorem}

\begin{theorem}[Power-of-three refinement]
For $N=3^b$, the same cardinality conclusion holds with
\[
 m\le(3X+1)2^{2^{15(3\cdot3^b)}}
       =(3X+1)2^{2^{45N}}.
\]
\end{theorem}

\begin{theorem}[Factor-sensitive refinement]
Let $N=2^\alpha3^bM$ with $\gcd(M,6)=1$, and suppose a certified parity word
and auxiliary integers satisfy the hypotheses of the formal factorized
transport theorem.  If $Z$ is the resulting effective scale, then the same
cardinality conclusion holds with
\[
 m\le2N(3X+1)2^{2^{15Z}}.
\]
The complete arithmetic certificate is stated in the Lean declaration
\texttt{direct\_to\_one\_logarithmic\_count\_factorized}.
\end{theorem}

\section{What the theorem means for Collatz dynamics}
The theorem gives a quantitative form of modular permeability.  No fixed
congruence class can be only sparsely accessible in the qualitative sense of
having a few exceptional predecessors: every class receives an explicitly
growing family of certified convergent integers.

This should not be confused with uniform randomness.  The result neither
compares different residue classes nor proves that their counts have an
asymptotic ratio of one.  The construction may be highly biased and its
numerical bounds are enormous.  What it rules out is a broad family of purely
modular obstructions to convergence-to-one membership.

The logarithmic denominator has a dynamical interpretation.  Effective
residue transport itself does not cause the loss.  The loss arises because
different fast targets can lie on the same orbit and therefore select the same
predecessor.  A Collatz orbit segment of logarithmic length can absorb at most
logarithmically many targets.  Thus the remaining step from the present
$X/\log X$ guarantee to a positive-density guarantee in each residue class is
a concrete multiplicity problem, not a modular reachability problem.

In a physical or statistical-mechanical analogy, the inverse dynamics has
full modular support but has not been shown to have a uniform stationary
distribution.  Coverage is analogous to irreducibility; positive density or
equidistribution would require quantitative control of mass concentration.
This analogy is interpretive only and is not used in the proof.

\section{Relation to the full Collatz conjecture}
The theorem does not prove that every integer reaches $1$, that almost every
integer in each residue reaches $1$, or even that the convergent integers have
positive lower density within each residue class.  Its exact finite lower
bound is of order $X/\log X$ in the target parameter, while transported
witnesses lie below a modulus-dependent constant multiple of $X$.

Nevertheless, the theorem sharpens the map of what a counterexample mechanism
would have to accomplish.  It cannot rely merely on excluding a residue class
from convergent inverse trees.  Any persistent nonconvergent component must
coexist with quantitatively many convergent starts in every congruence class.

\section{Further directions}
The most immediate problem is to replace the worst-case fibre bound $B+4$ by
a constant or a bounded average on a positive-density subfamily.  Such a result
would upgrade the construction toward positive density in every residue
class.  Computational experiments can measure fibre histograms for different
target selections and test whether collision multiplicity remains bounded,
has a light tail, or grows logarithmically.

Other directions include optimizing the enormous effective transport bounds,
comparing counts across residue classes, transporting stronger stopping-time
distributions, and replacing the external Effective Remainder Everywhere
dependency with an independent formalization.

\section{Lean formalization}
The development uses Lean 4.30.0-rc2 with a pinned Mathlib revision.  The
principal declarations are:
\begin{itemize}[leftmargin=2em]
\item \texttt{explicit\_logarithmic\_remainder\_transport\_ten46};
\item \texttt{direct\_to\_one\_logarithmic\_count\_ten46};
\item \texttt{direct\_to\_one\_count\_of\_selector};
\item \texttt{direct\_to\_one\_logarithmic\_count\_coprime\_six};
\item \texttt{direct\_to\_one\_logarithmic\_count\_pow\_three};
\item \texttt{direct\_to\_one\_logarithmic\_count\_factorized}.
\end{itemize}
For each, \texttt{\#print axioms} reports only \texttt{propext},
\texttt{Classical.choice}, and \texttt{Quot.sound}.  No \texttt{sorry} or
\texttt{admit} occurs in the audited theorem modules.

The initial public release does not redistribute Hallqvist's unlicensed Lean
source.  It records that source as an external prerequisite by repository URL,
filename, and hash.  Mazur's required source closure is redistributed under
its Apache-2.0 terms with its notice retained.

\section*{Authorship and AI disclosure}
The project architecture, selection of the synthesis problem, iterative
verification process, theorem packaging, and publication are attributed to
the named author.  OpenAI language models accessed through ChatGPT/Codex
contributed substantially to mathematical exploration, Lean proof generation,
debugging, and manuscript drafting.  The upstream Mazur development also
discloses substantial generative-AI assistance, and Lau's refinement discloses
autonomous AxiomProver contributions.  AI systems are not listed as authors or
independent referees.  The named author assumes responsibility for the claims
and released artifacts.

\begin{thebibliography}{99}
\bibitem{Everett1977}
C.~J. Everett,
\emph{Iteration of the number-theoretic function $f(2n)=n$, $f(2n+1)=3n+2$},
Advances in Mathematics 25 (1977), 42--45.

\bibitem{HallqvistERE}
A. Hallqvist,
\emph{Every Remainder Everywhere}, 2026,
Zenodo, DOI: \href{https://doi.org/10.5281/zenodo.23757322}{10.5281/zenodo.23757322}.

\bibitem{HallqvistEffective}
A. Hallqvist,
\emph{Effective Remainder Everywhere}, GitHub repository, 2026,
\url{https://github.com/Sio-MLST/Effective-Remainder-Everywhere}.

\bibitem{Lau2026}
K. Lau,
\emph{Explicit positive lower density for the Collatz Conjecture},
Hexagon 2610.00249, version 2, 2026.

\bibitem{MazurLog}
L. Mazur,
\emph{Explicit Positive-Density Collatz Convergence in Logarithmic Time},
manuscript version 2.1, 2026; formal source archived by ProofAtlas.

\bibitem{MazurPred}
L. Mazur,
\emph{Positive Lower Density of Collatz Predecessors},
manuscript, 2026; formal source archived by ProofAtlas.

\bibitem{Tao2022}
T. Tao,
\emph{Almost all orbits of the Collatz map attain almost bounded values},
Forum of Mathematics, Pi 10 (2022), e12,
DOI: \href{https://doi.org/10.1017/fmp.2022.8}{10.1017/fmp.2022.8}.

\bibitem{Terras1976}
R. Terras,
\emph{A stopping time on the positive integers},
Acta Arithmetica 30 (1976), 241--252.

\bibitem{Terras1979}
R. Terras,
\emph{On the existence of a density},
Acta Arithmetica 35 (1979), 101--102.

\bibitem{Wirsching1998}
G. J. Wirsching,
\emph{The Dynamical System Generated by the $3n+1$ Function},
Lecture Notes in Mathematics, Springer, 1998,
DOI: \href{https://doi.org/10.1007/BFb0095985}{10.1007/BFb0095985}.

\bibitem{Wirsching2003}
G. J. Wirsching,
\emph{On the problem of positive predecessor density in $3n+1$ dynamics},
Discrete and Continuous Dynamical Systems 9 (2003), 771--787,
DOI: \href{https://doi.org/10.3934/dcds.2003.9.771}{10.3934/dcds.2003.9.771}.
\end{thebibliography}

\end{document}
